\begin{figure}[htbp]
  \centering
  \resizebox{0.98\textwidth}{!}{%
  \begin{tikzpicture}[
    node distance=8mm and 12mm,
    layer/.style={rectangle, rounded corners, draw, fill=#1!15, minimum width=3.2cm, minimum height=1.0cm, align=center},
    note/.style={rectangle, draw, fill=#1!8, minimum width=5.2cm, minimum height=1.0cm, align=left},
    arr/.style={-Stealth, thick}
  ]
    \node[layer=blue] (l1) {第一层：教学级耦合波模型};
    \node[layer=teal, below=of l1] (l2) {第二层：单胞模筛选};
    \node[layer=orange, below=of l2] (l3) {第三层：有限阵列校核};
    \node[layer=red, below=of l3] (l4) {第四层：电热回写闭环};

    \node[note=blue, right=14mm of l1] (n1) {输入：$\kappa_1,\kappa_2,\gamma_i,\gamma_r$\\输出：模分裂趋势与排序直觉};
    \node[note=teal, right=14mm of l2] (n2) {输入：层栈与单胞几何\\输出：候选模与被动裕量};
    \node[note=orange, right=14mm of l3] (n3) {输入：孔径与边界假设\\输出：边缘泄漏、近远场稳定性};
    \node[note=red, right=14mm of l4] (n4) {输入：$J(\mathbf r),N(\mathbf r),T(\mathbf r)$\\输出：器件层面的通过/失败判定};

    \draw[arr] (l1) -- (l2);
    \draw[arr] (l2) -- (l3);
    \draw[arr] (l3) -- (l4);
    \draw[arr] (l1) -- (n1);
    \draw[arr] (l2) -- (n2);
    \draw[arr] (l3) -- (n3);
    \draw[arr] (l4) -- (n4);
  \end{tikzpicture}
  }
  \caption{算例桥接图：每一层都给出明确输入、输出与结论边界，读者可沿证据链追踪“本层到底证明了什么”。}
  \label{fig:ch23_two_layer_example_bridge}
\end{figure}
